% TOP_PROBLEM: {"id": 30006795, "title": "Sharp Erdos Distinct-Distances Conjecture in the Plane", "category_id": 6, "category": {"id": 6, "name": "geometry", "display_name": "Geometry", "description": "Euclidean and non-Euclidean geometry, geometric structures.", "slug": "geometry", "order_index": 6, "created_at": "2026-07-31T15:26:25.671Z"}, "status": "open"}
% FIRST_POSTED: 2026-09-27
% !TeX program = lualatex
% Compile twice: lualatex 30006795.tex
% All references and prose are embedded; no BibTeX database or external inputs.
\documentclass[11pt,a4paper]{article}
\usepackage[margin=24mm,headheight=14pt]{geometry}
\usepackage{fontspec}
\setmainfont{Latin Modern Roman}
\setsansfont{Latin Modern Sans}
\setmonofont{Latin Modern Mono}
\usepackage{amsmath,amssymb,mathtools}
\usepackage{unicode-math}
\setmathfont{Latin Modern Math}
\usepackage{microtype}
\usepackage{enumitem,xurl,needspace}
\usepackage[dvipsnames]{xcolor}
\usepackage{hyperref}
\hypersetup{unicode=true,colorlinks=true,linkcolor=MidnightBlue,urlcolor=MidnightBlue,
 pdftitle={Research Notebook: Section 30006795},pdfauthor={Research notebook},bookmarksnumbered=true}
\usepackage{bookmark}
\usepackage{tocloft}
\setlength{\cftsecnumwidth}{3em}
\cftsetpnumwidth{2.5em}
\cftsetrmarg{4em}
\usepackage{titlesec}
\titleformat{\section}{\Large\bfseries\raggedright}{\thesection}{0.7em}{}
\usepackage{fancyhdr}
\pagestyle{fancy}\fancyhf{}
\fancyhead[L]{\small\sffamily Open Problems | Research notebook}
\fancyhead[R]{\small 218 | 27 September 2026}
\fancyfoot[C]{\thepage}
\setlength{\parindent}{0pt}
\setlength{\parskip}{5pt plus 1pt}
\setlength{\emergencystretch}{3em}
\setcounter{tocdepth}{1}
\setcounter{secnumdepth}{1}
\setlist[itemize]{leftmargin=3em,itemsep=5pt,topsep=4pt,parsep=0pt}
\allowdisplaybreaks
\newcommand{\N}{\mathbb{N}}
\newcommand{\Z}{\mathbb{Z}}
\newcommand{\Q}{\mathbb{Q}}
\newcommand{\R}{\mathbb{R}}
\newcommand{\C}{\mathbb{C}}
\newcommand{\F}{\mathbb{F}}
\newcommand{\A}{\mathbb{A}}
\newcommand{\PP}{\mathbb P}
\newcommand{\E}{\mathbb{E}}
\newcommand{\eps}{\varepsilon}
\newcommand{\dd}{\,\mathrm d}
\newcommand{\sref}[2]{\hyperlink{source:#1:#2}{[S#2]}}
\newcommand{\eref}[2]{\hyperlink{extra:#1:#2}{[E#2]}}
% Turn the next switch off for a shorter reading copy. The quotations preserve
% every exactTarget field, including qualifications omitted from a short summary.
\newif\ifshowcatalogue
\showcataloguetrue
\newcommand{\cataloguescope}[1]{\ifshowcatalogue
 \paragraph{Exact scope in the supplied catalogue.}
 \begin{quote}\small #1\end{quote}\fi}

% Independently compilable extraction; original section text retained.
\numberwithin{equation}{section}
\begin{document}
\setcounter{section}{0}
% ================================================================
% problemId: 30006795
% Canonical problem ID: 30006795
% exactTarget SHA-256: ab249bc9dbc8d365dd25838bab9d0599e91bcb9be6771aa4a97805c5f8f683ad
\clearpage
\section*{\texorpdfstring{Sharp Erdos Distinct-Distances Conjecture in the Plane}{Sharp Erdos Distinct-Distances Conjecture in the Plane}}\label{30006795}
\subsection{Definitions and mathematical statement}
For a finite set $P\subset\R^2$ of distinct points, put
$D(P)=|\{|x-y|:x,y\in P,\ x\ne y\}|$, and let
$D(n)=\min_{|P|=n}D(P)$. The selected lower bound is
\[
 \exists c>0\ \exists n_0\ \forall n\ge n_0,
 \qquad D(n)\ge c\frac{n}{\sqrt{\log n}}.
\]
Together with the square-lattice comparison specified in the catalogue, this would determine the order up to constants. Distances are counted once regardless of multiplicity, and the metric is the usual Euclidean metric, not a generic norm. Improving a bound while leaving an extra power of $\log n$ does not reach this target.
\par\smallskip\textit{Formulation and concept references:} \sref{30006795}{1}.
\cataloguescope{Let D(n) be the minimum, over all sets of n distinct points in the Euclidean plane, of the number of distinct pairwise distances determined by the set. Prove D(n) = Omega(n/sqrt(log n)), matching up to a constant factor the square-lattice construction D(n) = O(n/sqrt(log n)).}
\subsection{Short English statement}
Must every n-point set in the Euclidean plane determine at least a constant times n divided by the square root of log n distinct distances?
% BEGIN RESEARCH ATTEMPT 218
\subsection{Research attempt}

\paragraph{Current frontier.}
The Guth--Katz theorem gives $D(P)\gg n/\log n$ for arbitrary Euclidean
point sets. The remaining square root of the logarithm is not a change
of constants. \eref{30006795}{1}, Theorem~1.1.
The catalogue's formulation reference studies \emph{typical norms}; its
statements about that class do not improve the Euclidean bound simply by
specializing the norm. \sref{30006795}{1}, Introduction. The source search for
this attempt did not identify a verified Euclidean theorem attaining the
selected bound. The argument below tests an energy improvement and then
replaces uniform pair-counting by a fractional allocation problem.
All research additions in this batch are dated 27 September 2026.

\paragraph{Attack route.}
Three routes are plausible: improve the equal-distance energy, discard
its highly concentrated fibers, or assign nonuniform weights to incidences
between points and their distances. The first route has a decisive lattice
obstruction, proved below. Merely truncating that energy still needs a
bound on the discarded mass. The third route has an exact max-flow
criterion; it avoids the false energy lemma but exposes a uniform subset
estimate that remains as difficult as the original assertion.

\paragraph{Attempt.}
\textbf{1. Why a stronger second-moment bound is the wrong goal.}
For a positive distance $t$ let
\[
 r_t=|\{(x,y)\in P^2:x\ne y,\ |x-y|=t\}|,
 \qquad S=\sum_t r_t=n(n-1),\qquad E(P)=\sum_t r_t^2.
\]
Cauchy--Schwarz gives $D(P)\ge S^2/E(P)$. Thus a bound
$E(P)\ll n^3\sqrt{\log n}$ would suffice, but it is false.
Here is a direct proof of the obstruction, independent of any numerical
experiment.

Take $P=\{1,\ldots,m\}^2$, $n=m^2$. A displacement
$h=(h_1,h_2)$ occurs as an ordered difference in
$(m-|h_1|)(m-|h_2|)$ pairs. In particular each displacement of Euclidean
norm at most $m/2$ has multiplicity at least $m^2/4$. Identify integer
vectors with Gaussian integers. For coprime positive integers $a,b$,
put $w=a+ib$ and, for $v\in\Z[i]\setminus\{0\}$, form
\[
 h=\overline w v,\qquad h'=wv.
\]
These are equal-length integer vectors. Restrict $a,b$ to $[R,2R]$ and
$|v|\le m/(8R)$. Then $|h|,|h'|<m/2$. There are at least $cR^2$
coprime pairs in that square for sufficiently large $R$, and at least
$c'(m/R)^2$ choices of $v$ when $R\le m/32$. For clarity, the former
count follows from
\[
 \sum_{q\le2R}\mu(q)
 \bigl(\lfloor2R/q\rfloor-\lfloor(R-1)/q\rfloor\bigr)^2
 =\frac{R^2}{\zeta(2)}+O(R\log R),
\]
using absolute convergence of $\sum q^{-2}$ and the elementary divisor
expansion for coprimality.

No two coprime positive $w$ with different slopes give the same ordered
pair $(h,h')$: the quotient $h'/h=w/\overline w$ determines the slope,
and primitivity determines $w$. For fixed $w$, $h$ determines $v$.
Use disjoint dyadic ranges for $a,b$, adjusting endpoints to avoid overlap.
Each range therefore contributes at least $cm^2$ different equal-length
displacement pairs, and there are $\gg\log m$ ranges. Each contributes
at least $m^4/16$ ordered point quadruples to $E(P)$. Consequently
\begin{equation}\label{30006795:lattice}
 E(P)\ge c m^6\log m\asymp n^3\log n.
\end{equation}
Thus even a successful proof of the requested distance bound cannot
supply the proposed universal second-moment improvement. The quadruples
counted here have nonzero displacements; including or excluding the zero
distance does not cause the obstruction.

\textbf{2. Fractional allocation in place of multiplicity.}
Let $T$ be the positive distance set, and join $p\in P$ to $t\in T$
when some $q\ne p$ satisfies $|p-q|=t$. Ask for nonnegative numbers
$w_{p,t}$ on these incidences such that
\begin{equation}\label{30006795:allocation}
 \sum_{t\sim p}w_{p,t}=1\quad(p\in P),\qquad
 \sum_{p\sim t}w_{p,t}\le B\quad(t\in T).
\end{equation}
Such an allocation immediately gives $n\le B|T|$. It can suppress large
pair multiplicities because a point contributes at most one unit of mass,
not one unit for each point at the same distance.

The exact existence criterion is
\begin{equation}\label{30006795:hall}
 |A|\le B\,|T(A)|\quad\hbox{for every }A\subseteq P,
 \qquad T(A)=\bigcup_{p\in A}\{t:t\sim p\}.
\end{equation}
To prove this, put capacities $1$ on source--point edges, $n+1$ on
point--distance edges and $B$ on distance--sink edges. A minimum cut of
capacity less than $n$ cannot cut a point--distance edge. If the retained
point subset is $A$, its optimal cut has capacity
$n-|A|+B|T(A)|$. The max-flow/min-cut theorem gives a flow of value $n$
exactly when \eqref{30006795:hall} holds. Necessity also follows directly by
summing \eqref{30006795:allocation}. Real capacities cause no integrality
issue; only a fractional allocation is requested.

A hoped-for geometric lemma would be \eqref{30006795:hall} with
$B=C\sqrt{\log(n+2)}$. It is tempting to regard this as much weaker than
the target, but that would be misleading. The case $A=P$ gives the target
immediately. Conversely the original bound, if true, would imply this
Hall estimate up to an absolute constant: for $|A|$ large,
\[
 |T(A)|\ge D(A)\ge c\frac{|A|}{\sqrt{\log |A|}}
            \ge c\frac{|A|}{\sqrt{\log(n+2)}}.
\]
For bounded nonempty $|A|$ and $n\ge2$, $T(A)$ is nonempty, and increasing
$C$ handles the remaining cases. Thus allocation is a useful exact
reformulation, not a secretly established intermediate theorem.

\textbf{3. A geometric test and an obstruction to naive truncation.}
A set with $m$ collinear points has at least $(m-1)/2$ distinct distances:
fix one of those points and note that a positive distance on its line has
at most two preimages. The same bound holds for $m$ points on a circle,
by intersecting that circle with a circle centered at a fixed point of it.
Consequently a line or circle carrying
$\gg n/\sqrt{\log n}$ points already proves the desired bound for that
configuration. A viable structural argument could therefore restrict
attention to sets with no such rich line or circle. This restriction is
only a solved subcase reduction, not an incidence theorem for the
remaining sets.

For a truncation level $L$, let $T_{\le L}=\{t:r_t\le L\}$. Then
\[
 D(P)\ge \frac{S-\sum_{r_t>L}r_t}{L}.
\]
To choose $L\asymp n\sqrt{\log n}$ and finish by this inequality one
must prove that the heavy fibers carry less than, say, $S/2$.
An energy bound $E\ll n^3\log n$ gives only
$\sum_{r_t>L}r_t\le E/L\ll n^2\sqrt{\log n}$, which is useless relative
to $S\asymp n^2$. This calculation identifies the missing tail estimate;
it does not justify it. The companion computation records exact grid
multiplicities, energy and distance counts at several sizes solely to
test these formulas.

\paragraph{Stress test.}
The lattice obstruction uses actual integer coordinates, not a norm
perturbation. The allocation criterion quantifies over \emph{all} point
subsets; checking only the full set assumes the target, and checking only
small sets cannot certify a flow. The rich-circle argument concerns
points of $P$ on the circle, not arbitrary distance circles with no such
points. No conclusion for all $n$ is drawn from a finite lattice scan.
Finally, the Guth--Katz energy method and the typical-norm result are
kept separate: neither supplies the unproved tail or Hall estimate here.

\paragraph{Outcome.}
\textbf{Outcome: Exploratory approach with unresolved gap.}

The attempt proves a sharp obstruction to improving the universal second
moment, an exact fractional-allocation reformulation, and elementary
rich-line/rich-circle subcases. These calculations explain why the
obvious logarithm-saving route fails; they are not a claimed improvement
of the best general lower bound or a claim of priority.
The unresolved task is a genuinely stronger geometric allocation or
tail estimate for arbitrary point configurations. No counterexample to
the selected distance conjecture has been found.
% END RESEARCH ATTEMPT 218
\subsection{Sources}
\begin{itemize}\raggedright
\item[\textnormal{[S1]}]\hypertarget{source:30006795:1}{} Unit and distinct distances in typical norms.
\url{https://arxiv.org/abs/2302.09058}.
{\small\textit{Catalogue formulation source.}}
% BEGIN ADDED REFERENCES 218
\item[\textnormal{[E1]}]\hypertarget{extra:30006795:1}{}
L. Guth and N. H. Katz, \emph{On the Erd\H{o}s distinct distances problem
in the plane}, Annals of Mathematics \textbf{181} (2015), 155--190,
Theorem~1.1 and the equal-distance-quadruple reduction in the Introduction.
\url{https://annals.math.princeton.edu/wp-content/uploads/annals-v181-n1-p02-p.pdf}.
Consulted 27 September 2026.
% END ADDED REFERENCES 218
\end{itemize}

\end{document}
